\begin{enumerate}
	\item Déterminons les coordonnées des vecteurs :

	$\vectt{AB} \begin{pmatrix}
		x_{\mathrm{B}} - x_{\mathrm{A}}\\
		y_{\mathrm{B}} - y_{\mathrm{A}}\\
		z_{\mathrm{B}} - z_{\mathrm{A}} \end{pmatrix} $
	\quad soit \quad $\vectt{AB} \begin{pmatrix} 3-2\\-2-1\\0-1 \end{pmatrix} $
	\quad finalement :\quad $\vectt{AB} \begin{pmatrix} 1\\-3\\-1\end{pmatrix} $.
	\quad De même, \quad $\vectt{AC} \begin{pmatrix} -2\\-2\\0 \end{pmatrix} $.

	Ces vecteurs ont des coordonnées non proportionnelles, donc ils ne sont pas colinéaires.

	Les points A, B et C ne sont pas alignés, donc ils définissent un plan.

	\item \begin{enumerate}
		\item On est dans un repère orthonormé, donc on peut utiliser la formule analytique du produit scalaire :

		\begin{itemize}
			\item $\vect{n} \cdot \vectt{AB} = 1 \times 1 -  1 \times  (-3) + 4  \times  (-1) = 1 + 3 - 4 = 0$;
			\item $\vect{n} \cdot \vectt{AC} = 1 \times (-2) -  1 \times  (-2) + 4  \times  0 = -2 + 2 + 0 = 0$;
			\item $\vectt{AB}$ et $\vectt{AC}$ sont deux vecteurs non colinéaires qui forment une base du plan (ABC).

			$\vect{n}$ est orthogonal à ces deux vecteurs, donc c'est un vecteur normal du plan (ABC).
		\end{itemize}
		\item Si $\vect{n}\begin{pmatrix} 1\\-1\\4 \end{pmatrix}$ est normal à (ABC), alors, par propriété, une équation de (ABC) peut avoir la forme :\quad $1\times x + (-1) \times y + 4 \times z + d = 0$, \quad où $d$ est un réel.

		Comme A est un point de (ABC), on doit avoir :\quad $x_{\mathrm{A}} - y_{\mathrm{A}} + 4z_{\mathrm{A}} + d =0$.

		$\begin{aligned}
			x_{\mathrm{A}} - y_{\mathrm{A}} + 4z_{\mathrm{A}} + d =0
			&\iff 2 - 1 + 4\times 1 + d= 0\\
			&\iff 5 + d = 0\\
			&\iff d = -5
		\end{aligned}		$

		Une équation cartésienne de (ABC) est donc bien :\quad $x-y+4z - 5 = 0$.
	\end{enumerate}

	\item Si la droite $\Delta$ est orthogonale au plan (ABC), alors, par définition, elle doit être dirigée par un vecteur normal au plan (ABC), donc ici, on en déduit que le vecteur $\vect{n} \begin{pmatrix} 1\\-1\\4 \end{pmatrix}$ dirige la droite $\Delta$.

	Comme la droite $\Delta$ passe par le point D$(0 ; 0 ; 2)$, une représentation paramétrique de la droite $\Delta$ est donnée par :\quad $\begin{cases}
		x=x_{\mathrm{D}} + x_{\vect{n}}\times t\\
		y=y_{\mathrm{D}} + y_{\vect{n}}\times t\\
		z=z_{\mathrm{D}} + z_{\vect{n}}\times t
	\end{cases}\text{avec } t \in \R$

	En remplaçant par les valeurs connues, on a bien :\quad $\begin{cases}
		x=t\\
		y=-t\\
		z=2+4t
	\end{cases}\text{avec } t \in \R$

	\item Le projeté orthogonal de D sur le plan (ABC) est l'intersection de la droite $\Delta$ et du plan (ABC).

	Comme c'est un point de la droite $\Delta$, il existe un réel $t_0$ tel que les coordonnées de ce point sont $(t_0 ; -t_0 ; 2+4t_0)$.

	C'est aussi un point de (ABC) sont ses coordonnées vérifient l'équation du plan.

	Autrement dit, $t_0$ est solution de l'équation :

	$\begin{aligned}[t]
		t_0 - (-t_0) + 4(2 + 4t_0) - 5 = 0
		&\iff t_0 + t_0 + 8 + 16t_0 - 5 = 0\\
		&\iff 18t_0 = -3\\
		&\iff t_0 = \dfrac{-1}{6}
	\end{aligned}$

	Le point H a donc pour coordonnées :\quad $\left(\dfrac{-1}{6} ; - \dfrac{-1}{6} ; 2 + 4\times \dfrac{-1}{6}\right)$\quad  soit $\mathrm{H}\left(-\dfrac{1}{6} ; \dfrac{1}{6} ; \dfrac{4}{3}\right) $.
	\item \begin{enumerate}
		\item Comme le repère est orthonormé :
		\begin{itemize}
			\item $\mathrm{AB} = \sqrt{(3 - 2)^2 + (-2 - 1)^2 + (0 - 1)^2} = \sqrt{1^2 + (-3)^2 + (-1)^2} = \sqrt{1 + 9 + 1} = \sqrt{11}$;
			\item $\mathrm{CB} = \sqrt{(3 - 0)^2 + (-2 - (-1))^2 + (0 - 1)^2} = \sqrt{3^2 + (-1)^2 + (-1)^2} = \sqrt{ 9 + 1 + 1} = \sqrt{11}$;
		\end{itemize}
		Les longueurs sont égales, donc le triangle est bien isocèle en B.

		\item \begin{itemize}
			\item Dans un triangle isocèle, la hauteur issue du sommet principal est aussi une médiane (et une médiatrice).

			La hauteur issue de B est donc le segment [BI], où I est le milieu de [AC].

			I a donc pour coordonnées :\quad I$\left(\dfrac{x_{\mathrm{A}} + x_{\mathrm{C}}}{2} ; \dfrac{y_{\mathrm{A}} + y_{\mathrm{C}}}{2} ; \dfrac{z_{\mathrm{A}} + z_{\mathrm{C}}}{2}\right) $\quad soit \quad $(1 ; 0 ; 1)$.

			\item L'aire d'un triangle est donnée par :\quad $\mathcal{A} = \dfrac{\text{base} \times \text{hauteur}
			}{2}$.

			Ici, on choisit comme base le segment [AC] et la hauteur correspondante, [BI].

			\item La distance AC est :\quad $\mathrm{AC} = \sqrt{(0 - 2)^2 + (-1 - 1)^2 + (1 - 1)^2} = \sqrt{8} = 2\sqrt{2}$.

			La distance BI est :\quad $\mathrm{BI} = \sqrt{(3 - 1)^2 + (-2 - 0)^2 + (0 - 1)^2} = \sqrt{9} = 3$.
		\end{itemize}
		On a donc : \quad $\mathcal{A}_{\mathrm{ABC}} = \dfrac{\mathrm{AC} \times \mathrm{BI}}{2} = \dfrac{2\sqrt{2} \times 3}{2} = 3\sqrt{2}$.
	\end{enumerate}

	\item \begin{enumerate}
		\item Avec les questions précédentes, on a calculé l'aire du triangle ABC, qui est l'une des quatre bases possibles pour le tétraèdre.

		Si on choisit ABC comme base de ABCD, alors la hauteur correspondante est la distance séparant le quatrième sommet (ici D), du plan contenant la base ABC, c'est-à-dire le plan (ABC).

		Cette distance est donc la distance séparant D de son projeté orthogonal sur (ABC), c'est-à-dire H, dont on a déterminé les coordonnées à la question \textbf{4.}.

		$\mathrm{DH} = \sqrt{\left(\dfrac{-1}{6}\right)^2 + \left(\dfrac{1}{6}\right)^2 + \left(\dfrac{4}{3} - 2\right) ^2 } = \sqrt{\dfrac{18}{36}} = \sqrt{\dfrac{1}{2}} = \dfrac{\sqrt{2}}{2}$.

		On a donc : \quad $\mathcal{V}_{\mathrm{ABCD}} = \dfrac{1}{3} \times \mathcal{A}_{\mathrm{ABC}} \times \mathrm{DH} = \dfrac{1}{3} \times 3\sqrt{2} \times \dfrac{\sqrt{2}}{2} = \dfrac{1}{3} \times \dfrac{3 \times 2}{2} = 1$.

		Le volume du tétraèdre est bien égal à 1.

		\item En faisant le même raisonnement, mais avec la base BCD et la hauteur associée, issue de A, que l'on peut noter $h$, on en déduit :\quad $\mathcal{V}_{\mathrm{ABCD}} = \dfrac{1}{3} \times \mathcal{A}_{\mathrm{BCD}} \times h$.

		En remplaçant par les valeurs connues, on a :\quad $1 = \dfrac{1}{3} \times \mathcal{A}_{\mathrm{BCD}} \sqrt{2} \iff \mathcal{A}_{\mathrm{BCD}} = \dfrac{1 \times 3}{\sqrt{2}} = \dfrac{3\sqrt{2}}{2}$.

		L'aire du triangle BCD est donc égale à $\dfrac{3\sqrt{2}}{2}$.
	\end{enumerate}

	\item \begin{enumerate}
		\item Avec $\mathrm{D}_k$ ayant pour coordonnées $(0 ; 0 ; k)$ (le point D du début de l'énoncé est donc $\mathrm{D}_2$).

		On va chercher quand ces coordonnées vérifient l'équation du plan (ABC).

		$\begin{aligned}
			\mathrm{D}_k \in (\mathrm{ABC})
				&\iff x_{\mathrm{D}_k} - y_{\mathrm{D}_k} + 4z_{\mathrm{D}_k} - 5 = 0\\
				&\iff 0 - 0 +4k  - 5 = 0\\
				&\iff k = \dfrac{5}{4}
 		\end{aligned}$

 		Seul le point $\mathrm{D}_{\frac54}$ est sur le plan (ABC).

 		Son projeté orthogonal sur (ABC) est alors lui-même, par  définition du projeté orthogonal.

		\item Déterminons les coordonnées du vecteur $\vect{\mathrm{AD}_k}$ : \quad $\vect{\mathrm{AD}_k} \begin{pmatrix} -2\\-1\\k-1 \end{pmatrix} $.\quad On rappelle : \quad $\vect{n} \begin{pmatrix} 1\\-1\\4 \end{pmatrix}$.

		Quelle que soit la valeur de $k$, ces vecteurs ne sont pas colinéaires : ils ont des ordonnées égales, mais des abscisses différentes.

		Donc le vecteur $\vect{\mathrm{AD}_k}$ n'est jamais orthogonal au plan (ABC) et donc A n'est jamais le projeté orthogonal de $\mathrm{D}_k$ sur (ABC).
	\end{enumerate}
\end{enumerate}
